\documentclass[11pt,a4paper]{article}
\usepackage[utf8]{inputenc}
\usepackage[T1]{fontenc}
\usepackage[brazil]{babel}
\usepackage[a4paper,landscape,margin=1.5cm]{geometry}
\usepackage{array}
\usepackage{booktabs}
\usepackage{longtable}
\usepackage{ragged2e}
\setlength{\LTleft}{0pt}
\setlength{\LTright}{0pt}
\renewcommand{\arraystretch}{1.18}
\newcolumntype{P}[1]{>{\RaggedRight\arraybackslash}p{#1}}
\begin{document}
\begin{center}
{\Large\bfseries Cronograma para redação de dissertação}\par
\vspace{0.2em}
{\large Técnicas de avaliação de sistemas RAG: aplicação empírica a discursos parlamentares do Senado Federal}\par
\vspace{0.3em}
Fabricio Fernandes Santana\\
Orientador: Prof. Dr. Marcelo Pita
\end{center}
\vspace{0.5em}
\small
\begin{longtable}{P{2.7cm}P{6.0cm}P{6.0cm}P{5.0cm}}
\caption{Cronograma de atividades, entregas e resultados esperados}\label{tab:cronograma}\\
\toprule
\textbf{Período} & \textbf{Atividade} & \textbf{Entrega para o orientador} & \textbf{Resultado esperado} \\
\midrule
\endfirsthead
\toprule
\textbf{Período} & \textbf{Atividade} & \textbf{Entrega para o orientador} & \textbf{Resultado esperado} \\
\midrule
\endhead
\bottomrule
\endlastfoot
01/08/2026 & Alinhar expectativas com o orientador e validar pergunta de pesquisa, objetivos, hipóteses, escopo e coleta de dados & Pré-projeto revisado; pauta de decisões metodológicas & Escopo da qualificação confirmado \\
01/09/2026 & Consolidar o referencial teórico e a metodologia para qualificação & Capítulos de introdução, referencial teórico, hipóteses e metodologia em versão expandida & Estrutura da qualificação validada \\
01/10/2026 & Fechar o texto para qualificação e preparar a apresentação & Versão quase final da qualificação; slides de apresentação de até 15 minutos & Aprovação do orientador para envio à banca \\
até 07/11/2026 & Enviar formalmente o material à banca de qualificação & Documento de qualificação enviado à banca & Cumprimento da antecedência mínima de 10 dias \\
17/11/2026 & Realizar a qualificação da dissertação & Apresentação oral e defesa do projeto perante a banca & Recebimento das recomendações da banca \\
18/11--15/12/2026 & Organizar e incorporar inicialmente as recomendações da banca & Matriz pós-qualificação com ajustes obrigatórios, recomendados e opcionais & Plano de ajustes pós-qualificação validado \\
16/12/2026--31/01/2027 & Preparar a execução empírica durante a licença & Corpus localizado; ambiente técnico preparado; checklist de execução; método ajustado após a banca & Condições técnicas e metodológicas preparadas para início da licença \\
01/02--14/02/2027 & Planejar o início do afastamento e organizar os instrumentos de acompanhamento & Plano de trabalho do primeiro período da licença; agenda de reuniões; repositório organizado & Prioridades da licença definidas com o orientador \\
15/02--21/02/2027 & Aprofundar a revisão bibliográfica e documental dirigida e realizar ajustes finais pós-qualificação & Matriz teórica revisada; lista final de referências essenciais & Aprofundamento da revisão bibliográfica e documental \\
22/02--28/02/2027 & Consolidar o referencial teórico e o percurso metodológico & Capítulos de introdução, referencial, hipóteses e metodologia revisados & Consolidação do referencial teórico e metodológico \\
01/03--07/03/2027 & Congelar o corpus e preparar tecnicamente o pipeline & Manifesto de dados; regras de inclusão e exclusão; parâmetros preliminares do pipeline RAG & Percurso metodológico consolidado e corpus versionado \\
08/03--14/03/2027 & Construir a bateria de perguntas e o conjunto de referência & Bateria de perguntas por categoria; evidências documentais esperadas & Projeto analítico consolidado \\
15/03--19/03/2027 & Validar o material produzido no primeiro período da licença & Sumário estruturado da dissertação; capítulos iniciais revisados; protocolo avaliativo v1 & Versão preliminar dos capítulos iniciais e da metodologia \\
20/03--31/03/2027 & Revisar, com o orientador, o material produzido no primeiro período da licença & Devolutiva sobre corpus, perguntas, protocolo e capítulos & Ajustes críticos identificados antes da execução empírica \\
01/04--04/04/2027 & Preparar a execução empírica final & Lista fechada de ajustes antes da execução do protocolo & Protocolo pronto para aplicação \\
05/04--11/04/2027 & Executar o pipeline RAG e coletar os resultados & Resultados brutos: evidências recuperadas e respostas geradas & Execução de testes, experimentos e análises preliminares \\
12/04--18/04/2027 & Realizar a avaliação humana e a avaliação baseada no paradigma LLM-as-a-judge & Planilha de julgamentos humanos; resultados automatizados; registro dos parâmetros de execução & Sistematização dos resultados parciais \\
19/04--23/04/2027 & Analisar os resultados e redigir os capítulos empíricos & Capítulo de resultados v1; análise preliminar das hipóteses H1, H2 e H3 & Minuta dos capítulos de desenvolvimento e análise inicial dos resultados \\
24/04--26/04/2027 & Revisar a versão final para a banca de defesa & Versão completa da dissertação para revisão conclusiva do orientador & Autorização para depósito da versão de defesa \\
até 27/04/2027 & Depositar ou enviar a versão à banca de defesa & Dissertação enviada à banca & Cumprimento da antecedência mínima de 10 dias para a defesa \\
28/04--30/04/2027 & Preparar a apresentação oral da defesa & Slides; roteiro de apresentação de 20 minutos; simulação de arguição & Defesa preparada \\
07/05/2027 & Defender a dissertação & Apresentação oral, arguição e deliberação da banca & Aprovação, aprovação com ajustes ou reformulação \\
08/05--15/05/2027 & Organizar as recomendações da banca & Matriz de correções pós-defesa & Plano de ajustes pós-defesa definido \\
16/05--30/05/2027 & Incorporar os ajustes iniciais recomendados pela banca & Versão pós-defesa v1 para o orientador & Material preparado para consolidação final durante o terceiro período da licença \\
31/05--06/06/2027 & Incorporar as recomendações da banca aos capítulos afetados & Versão revisada dos capítulos afetados & Análise final dos resultados e discussão incorporadas \\
07/06--13/06/2027 & Consolidar a conclusão, as recomendações e os limites da pesquisa & Conclusão revisada; recomendações de uso cauteloso; limites da pesquisa & Capítulos finais consolidados \\
14/06--20/06/2027 & Revisar o texto do ponto de vista acadêmico, metodológico e bibliográfico & Versão completa revisada & Texto revisado do ponto de vista acadêmico, metodológico e bibliográfico \\
21/06--25/06/2027 & Normalizar a dissertação conforme normas acadêmicas & Versão normalizada da dissertação & Adequação às normas acadêmicas e formatação final \\
26/06--30/06/2027 & Fechar a versão final e preparar o depósito institucional & Versão final depositável; plano do artigo derivado & Dissertação concluída em versão final, texto revisado e normalizado, versão pronta para depósito institucional \\
\end{longtable}
\end{document}
